\documentclass[12pt]{article}
\usepackage{report}

\begin{document}

\reporttitle{CS340400 Compiler Design - Homework 3 Implementation Report}

\section{Overview}

This assignment extends the HW2 parser into a compiler: it reads a subset of C and writes RISC-V
assembly to \file{codegen.S}. For testing, the TA's \file{main.c} calls the generated \code{codegen()},
which in turn calls the Arduino functions \code{digitalWrite()} and \code{delay()}; running on the Spike
simulator prints a log of those calls.

\begin{center}
\begin{tabularx}{0.95\linewidth}{@{}lL@{}}
\toprule
Files & \file{scanner.l}, \file{parser.y}, \file{makefile} \\
Grammar & byacc: \textbf{0 conflicts} \\
Build & GNU gcc 13.3, \textbf{0 warnings} (also with \code{-Wall -Wextra}) \\
Public tests & the TA's 9 test programs on Spike --- \textbf{9/9 pass} \\
Own tests & 6 programs in \file{tests/extra/} standing in for the hidden tests --- \textbf{6/6 pass} \\
RV32 & the same compiler with \code{--rv32}, all 15 programs on RV32 Spike --- \textbf{15/15 pass} \\
Usage & \code{make codegen && ./codegen < input.c}; the output goes to \file{codegen.S} \\
\bottomrule
\end{tabularx}
\end{center}

The design goals are a clear separation between parsing and code generation, correct recursive
functions, and a proper stack frame that follows the RISC-V calling convention.

\subsection{Constraints from the specification}

\begin{itemize}
  \item Only \file{scanner.l}, \file{parser.y} and \file{makefile} are submitted, so the AST, the symbol
        table and the code generator all live in \file{parser.y}.
  \item \textbf{No static optimisation}: no constant folding, constant propagation or dead-code elimination.
  \item RV64: 8-byte memory accesses with \code{ld} and \code{sd}.
  \item Calls follow the RISC-V calling convention: arguments in \code{a0}--\code{a7}, the result in \code{a0}.
  \item \code{HIGH} equals 1 and \code{LOW} equals 0.
\end{itemize}

\section{Architecture}

\begin{center}
\begin{tikzpicture}[
    box/.style={draw=headblue, rounded corners=2pt, fill=codebg, minimum height=8mm,
                inner xsep=5pt, font=\small\headfont, align=center},
    lab/.style={font=\scriptsize\headfont, text=gray, align=center},
    >={Stealth[length=2mm]}]
  \node[box] (scan) {\file{scanner.l}};
  \node[box, right=9mm of scan] (gram) {grammar\\[-2pt]{\scriptsize same as HW2}};
  \node[box, right=9mm of gram] (ast) {AST};
  \node[box, right=9mm of ast] (gen) {codegen walk\\[-2pt]{\scriptsize + symbol table}};
  \node[box, right=9mm of gen] (out) {\file{codegen.S}};
  \draw[->] (scan) -- (gram);
  \draw[->] (gram) -- (ast);
  \draw[->] (ast) -- (gen);
  \draw[->] (gen) -- (out);
  \draw[headblue, dashed, rounded corners=3pt]
     ($(gram.north west)+(-2mm,5mm)$) rectangle ($(gen.south east)+(2mm,-3mm)$);
  \node[lab, anchor=south] at ($(ast.north)+(0,5mm)$) {\file{parser.y}};
\end{tikzpicture}
\end{center}

\file{parser.y} is divided into five numbered blocks, each introduced by a banner comment:

\begin{center}
\begin{tabularx}{0.95\linewidth}{@{}lL@{}}
\toprule
Block & Contents \\
\midrule
{[1]} AST & node kinds, the \code{Node} struct, constructors (in the prologue, used by the actions) \\
{[2]} Grammar & identical to HW2; the actions build AST nodes instead of strings \\
{[3]} Symbol table & a scope stack stored in an array \\
{[4]} Code generation & expressions, statements, functions, globals \\
{[5]} \code{main} & command-line options, target selection, opening \file{codegen.S} \\
\bottomrule
\end{tabularx}
\end{center}

\subsection{Why build an AST first}

The specification's hints and the lecture slides emit assembly directly from the yacc actions with
\code{fprintf}. I build an AST first and generate code in a separate walk after parsing, because:

\begin{enumerate}
  \item \textbf{The frame size must be known up front.} The prologue allocates the whole stack frame at
        once, but the number of locals is known only after the whole body is parsed. With an AST, the
        generator counts them first (\code{count_slots}) and then emits the prologue.
  \item \textbf{Labels are easier to manage.} \code{break} and \code{continue} jump to labels outside the
        loop; a walk over the AST tracks the enclosing loops with a simple stack.
  \item \textbf{Pointer arithmetic needs types.} \code{p + 1} must be scaled when \code{p} is a pointer; the
        walk can consult the symbol table to decide.
  \item \textbf{Switching targets is cheap.} RV32 and RV64 differ in a handful of parameters, which the walk
        swaps without touching the grammar.
\end{enumerate}

\section{Scanner}

The scanner is HW2's with three extra rules. Two of them turn the Arduino constants into integer
literals, so neither the parser nor the generator needs to know about them:

\snippet[language=flex]{../scanner.l}{"HIGH"}{+2}

The third rule, for \code{uint32_t}, belongs to the Bonus and is described in the Bonus report.

\section{AST}

\subsection{The node structure}

\snippet[language=yacc]{../parser.y}{typedef struct Node}{@end}

Every node uses the same struct, and \code{kind} says what it is. A node has at most four children
(\code{a}--\code{d}); \code{for} uses all four: initialiser, condition, step and body. Lists of
parameters, arguments or statements are chained through \code{next}. \code{line} records the line number
for error messages.

\begin{center}
\begin{tabularx}{\linewidth}{@{}llL@{}}
\toprule
Group & Kind & Fields used \\
\midrule
Expression & \code{E_NUM} & \code{ival}: integer and character literals (\code{'a'} becomes 97 while building the AST) \\
 & \code{E_VAR} & \code{name} \\
 & \code{E_BIN}, \code{E_ASSIGN} & \code{op}, \code{a} (left), \code{b} (right) \\
 & \code{E_PRE}, \code{E_POST}, \code{E_UNARY} & \code{op}, \code{a} \\
 & \code{E_DEREF}, \code{E_ADDR}, \code{E_CAST} & \code{a} (\code{E_CAST} also sets \code{is_ptr}, \code{is_unsigned}) \\
 & \code{E_CALL}, \code{E_INDEX} & \code{name}, \code{a} (arguments / index) \\
Statement & \code{S_IF} & \code{a} (condition), \code{b} (then), \code{c} (else, may be NULL) \\
 & \code{S_WHILE}, \code{S_DO} & condition and body \\
 & \code{S_FOR} & \code{a}, \code{b}, \code{c}, \code{d} \\
 & \code{S_SWITCH}, \code{S_CASE} & a \code{S_CASE} with \code{a == NULL} is \code{default} \\
 & \code{S_BLOCK}, \code{S_DECLS} & \code{a} (statement list, declarator list) \\
Declaration & \code{D_VAR} & \code{name}, \code{is_ptr}, \code{a} (array dimensions), \code{b} (initialiser) \\
 & \code{D_FUNC} & \code{name}, \code{is_ptr}, \code{a} (parameters), \code{b} (body, NULL for a prototype) \\
\bottomrule
\end{tabularx}
\end{center}

\subsection{Grammar actions}

The grammar is the same as HW2's; only the actions change, from building strings to building nodes:

\snippetafter{start program}
\snippet[language=yacc]{../parser.y}{^expr}{+3}
\snippet[language=yacc]{../parser.y}{'(' expr ')'}{+3}

Parentheses only affect the shape of the tree, so \code{'(' expr ')'} simply returns the inner node.

Types are simplified as well. The generator does not distinguish \code{int}, \code{char} and \code{long}:
every integer and pointer occupies one machine word. The semantic value of \code{type_spec} is therefore no
longer a string but a single flag, whether the type is unsigned (used by the Bonus). Whether a variable is
a pointer comes from the \code{*} in its declarator and is stored in \code{is_ptr}.

\section{Symbol Table}

\snippetafter{SYM_LOCAL, SYM_GLOBAL}
\snippet[language=yacc]{../parser.y}{typedef struct}{@end}
\snippet[language=yacc]{../parser.y}{static void scope_push}{@end}
\snippet[language=yacc]{../parser.y}{static Sym *sym_lookup}{@end}

The symbol table is an array used as a stack:

\begin{itemize}
  \item Every symbol records the scope depth it was declared at. Entering \code{\{ \}} increments
        \code{depth}; leaving it pops every symbol of that depth off the top.
  \item Lookup scans \emph{from the top down}, so the first match is the innermost declaration and
        shadowing works without extra effort.
  \item There are three kinds of symbols: \code{SYM_LOCAL} (locals and parameters, with an offset from
        \code{s0}), \code{SYM_GLOBAL} (addressed with \code{la}) and \code{SYM_FUNC} (records whether the
        function returns a pointer).
\end{itemize}

Because the generator walks the AST in order and adds a symbol when it reaches the declaration, a name
is visible from its declaration onwards, as in C.

\section{Stack Frame}

\subsection{Layout}

\begin{center}
\begin{tikzpicture}[
    cell/.style={draw=headblue, fill=codebg, minimum width=48mm, minimum height=6.5mm,
                 font=\small\headfont, inner sep=0pt},
    tmp/.style={cell, fill=white, dashed},
    off/.style={font=\scriptsize\ttfamily, text=gray, anchor=west},
    note/.style={font=\scriptsize\headfont, text=gray, anchor=east, align=right},
    >={Stealth[length=2mm]}]
  \node[cell] (ra) {\code{ra}};
  \node[cell, below=0pt of ra] (s0) {old \code{s0}};
  \node[cell, below=0pt of s0] (p1) {parameter 1 (from \code{a0})};
  \node[cell, below=0pt of p1] (p2) {parameter 2 (from \code{a1})};
  \node[cell, below=0pt of p2, minimum height=9mm] (loc) {locals, one slot per declaration};
  \node[cell, below=0pt of loc, minimum height=9mm] (arr) {array \code{arr[n]}};
  \node[tmp, below=0pt of arr, minimum height=9mm] (t) {temporaries (push/pop)};
  \node[off] at ($(ra.east)+(2mm,0)$) {s0 - 8};
  \node[off] at ($(s0.east)+(2mm,0)$) {s0 - 16};
  \node[off] at ($(p1.east)+(2mm,0)$) {s0 - 24};
  \node[off] at ($(p2.east)+(2mm,0)$) {s0 - 32};
  \node[off] at ($(arr.east)+(2mm,-2.2mm)$) {arr[0]: lowest address};
  \draw[->, gray] ($(ra.north west)+(-14mm,0)$) node[note, left] {\code{s0} = caller's \code{sp}} -- (ra.north west);
  \draw[->, gray] ($(arr.south west)+(-14mm,0)$) node[note, left] {\code{sp} = \code{s0} - frame\\(16-byte aligned)} -- (arr.south west);
  \draw[->, gray] ($(ra.north east)+(44mm,0)$) -- node[font=\scriptsize\headfont, right, text=gray] {lower addresses} ($(t.south east)+(44mm,0)$);
\end{tikzpicture}
\end{center}

\begin{itemize}
  \item \textbf{Every integer and pointer takes one word} (8 bytes on RV64). Array elements are one word
        apart and pointer arithmetic scales by one word, so no per-type sizes are needed.
  \item \textbf{Every declaration gets its own slot, and slots are never reused}, even for variables in
        blocks that are never live at the same time. The frame is larger, but every offset is fixed when
        the code is generated, which leaves little room for mistakes.
  \item \textbf{Parameters are stored in slots too.} Right after the prologue, \code{a0}--\code{a7} are saved
        to the frame. Parameters and locals are then handled identically, and \code{&param} has an address:
        the public test Function/0 calls \code{modulo(&a, &b)} on \code{gcd}'s parameters.
  \item \textbf{Element 0 of an array is at the lowest address}, so \code{arr + i} moves towards higher
        addresses as C requires.
\end{itemize}

\subsection{Sizing the frame before emitting it}

\snippet[language=yacc]{../parser.y}{static int count_slots}{@end}

\code{count_slots} walks the function body and adds up the slots its declarations need: \code{n} for an
array, 1 for anything else, plus one per \code{switch} to hold the value being matched. The frame is then
two slots for \code{ra} and the old \code{s0}, plus parameters and locals, rounded up to a multiple of 16
because the ABI keeps \code{sp} 16-byte aligned:

\snippet[language=yacc]{../parser.y}{int frame = 2 * W}{+1}

\subsection{Prologue and epilogue}

\snippet[language=yacc]{../parser.y}{static void gen_function}{@end}

For \code{int add(int a, int b) \{ int c = a + b; return c; \}} the generator produces:

\begin{lstlisting}[language=riscv,title={\codetitle{codegen.S (RV64), comments added}}]
.global add
add:
  addi sp, sp, -48          # prologue: 2 + 2 + 1 = 5 slots = 40 bytes, rounded up to 48
  sd ra, 40(sp)
  sd s0, 32(sp)
  addi s0, sp, 48
  sd a0, -24(s0)            # parameters go to their own slots
  sd a1, -32(s0)
  addi t0, s0, -24          # load a
  ld t0, 0(t0)
  addi sp, sp, -16          # push a
  sd t0, 0(sp)
  addi t0, s0, -32          # load b
  ld t0, 0(t0)
  ld t1, 0(sp)              # pop a into t1
  addi sp, sp, 16
  add t0, t1, t0            # a + b
  sd t0, -40(s0)            # store into c
  addi t0, s0, -40          # return c
  ld t0, 0(t0)
  mv a0, t0
  j .L1
.L1:
  ld ra, -8(s0)             # epilogue
  mv t2, s0
  ld s0, -16(t2)
  mv sp, t2
  jalr zero, 0(ra)
\end{lstlisting}

The epilogue restores \code{sp} from \code{s0} instead of adding the frame size back. This stays correct
even if \code{return} is reached while temporaries are still pushed on the stack. Every \code{return}
jumps to the single epilogue label. When the frame exceeds 2047 bytes, a 12-bit immediate no longer fits,
and \code{addi_any} and \code{frame_mem} fall back to \code{li t2, imm} followed by \code{add}.

\section{Expressions}

\subsection{A stack machine}

Expressions are evaluated like a stack machine, with two rules:

\begin{enumerate}
  \item Every expression leaves its result in \code{t0}.
  \item A binary operation evaluates its left side, pushes \code{t0}, evaluates its right side, and pops the
        left value into \code{t1}. At that point \code{t1} holds the left operand and \code{t0} the right one.
\end{enumerate}

\snippet[language=yacc]{../parser.y}{static void push(}{2@end}

Each push takes 16 bytes (8 are used) so that \code{sp} stays 16-byte aligned, even when a call happens in
the middle of an expression.

The code is long, but every expression is handled the same way, which makes it easy to verify. Since
temporaries live on the stack rather than in registers, \textbf{nothing has to be saved around a call}:
no caller-saved register ever holds a live value. The specification forbids optimisation anyway, so code
size is not a goal.

\subsection{Binary operators}

\snippet[language=yacc]{../parser.y}{static void gen_binary}{@end}

After the push and pop, \code{gen_binary} emits the instruction for the operator:

\begin{center}
\begin{tabularx}{\linewidth}{@{}lp{6.2cm}L@{}}
\toprule
Operator & Instructions (\code{t1} left, \code{t0} right) & Note \\
\midrule
\code{+ - *} & \code{add}, \code{sub}, \code{mul} & pointer operands are scaled, see Section~\ref{sec:ptr} \\
\code{/ \%} & \code{div}, \code{rem} & RISC-V truncates towards zero, as C does \\
\code{& | ^} & \code{and}, \code{or}, \code{xor} & \\
\code{<< >>} & \code{sll}, \code{sra} & \code{>>} on a signed value is an arithmetic shift \\
\code{<} & \code{slt t0, t1, t0} & \\
\code{>} & \code{slt t0, t0, t1} & operands swapped \\
\code{<=} & \code{slt t0, t0, t1}; \code{xori t0, t0, 1} & \code{a <= b} is \code{!(b < a)} \\
\code{>=} & \code{slt t0, t1, t0}; \code{xori t0, t0, 1} & \code{a >= b} is \code{!(a < b)} \\
\code{==} & \code{sub}; \code{seqz} & equal when the difference is 0 \\
\code{!=} & \code{sub}; \code{snez} & \\
\bottomrule
\end{tabularx}
\end{center}

Every comparison yields 0 or 1, as in C. The \code{uns} flag and \code{sltu}/\code{divu}/\code{remu}/\code{srl}
are the Bonus's unsigned variants.

\subsection{Short-circuit evaluation}

\code{&&} and \code{||} must not evaluate both sides first: in \code{a && f()}, \code{f} must not be called
when \code{a} is 0. They are therefore handled at the top of \code{gen_binary}. After the left side, a 0 for
\code{&&} or a non-zero value for \code{||} jumps straight to \code{short_label}, which sets the result to 0 or 1
respectively. Otherwise the right side is evaluated and normalised to 0 or 1 with \code{snez}. The test
\file{short_circuit} checks that a right-hand call with a side effect is skipped.

\subsection{Lvalues and assignment}

Everything that needs an \emph{address} --- the left side of an assignment, \code{&}, \code{++} --- calls
\code{gen_addr}, which leaves the address in \code{t0}:

\snippet[language=yacc]{../parser.y}{static void gen_addr}{@end}

\begin{center}
\begin{tabularx}{\linewidth}{@{}lL@{}}
\toprule
Node & Address \\
\midrule
\code{E_VAR} & a local is at \code{s0 + offset}; a global is \code{la t0, name} \\
\code{E_DEREF} (\code{*p}) & the value of \code{p} is the address, so just \code{gen_expr(p)} \\
\code{E_INDEX} (\code{a[i]}) & base plus \code{i << W_SHIFT}; the base of an array is its own address, a pointer is loaded first \\
\bottomrule
\end{tabularx}
\end{center}

Reading a variable is \code{gen_addr} followed by a load. The one exception is the \textbf{name of an array}:
its value is its address, so no load follows. This is what makes \code{int *p = arr;} and \code{*(arr + i)} work:

\snippet[language=yacc]{../parser.y}{if (!s->is_array) ins(}{+1}

An assignment pushes the address of the left side, evaluates the right side, pops the address into
\code{t1} and stores with \code{sd t0, 0(t1)}. The value stays in \code{t0}, so chained assignments such as
\code{a = b = 1} work. Prefix and postfix \code{++}/\code{--} differ only in which value is left in \code{t0}:

\snippet[language=yacc]{../parser.y}{int step = }{break;}

\subsection{Pointer arithmetic}
\label{sec:ptr}

In C, pointer arithmetic counts elements: \code{p + 1} points to the next element, so the byte offset is
scaled by the element size. \code{is_ptr_expr} decides whether an expression yields a pointer:

\snippet[language=yacc]{../parser.y}{static int is_ptr_expr}{@end}

\code{gen_binary} uses the result for both operands: for \code{+}, the non-pointer operand is shifted left by
\code{W_SHIFT} (3 on RV64); for \code{-}, a pointer minus an integer scales the integer, and a pointer minus a
pointer divides the byte difference with \code{srai} to get an element count. The public test Jump/1 needs
this for \code{int *f1 = fib + i - 1}: \code{fib + i} is a pointer, so subtracting 1 is scaled too. For the same
reason \code{p++} steps by one word rather than by 1.

\section{Control Flow}

\subsection{Labels and conditional jumps}

Each control structure takes fresh numbers from \code{new_label()} and emits them as \code{.L<n>}. A false
condition has to jump away, but a RISC-V conditional branch such as \code{beqz} reaches only $\pm$4\,KB, which
a long loop body can exceed. The generator therefore always branches \emph{over} an unconditional jump:

\snippet[language=yacc]{../parser.y}{static void jump_if_zero}{@end}

\code{j} reaches $\pm$1\,MB, far more than any program in this assignment needs.

\subsection{Loops}

The labels are arranged as follows; \code{break} jumps to \code{end}, and \code{continue} to the label marked
$\star$:

\begin{center}
\begin{tabularx}{\linewidth}{@{}lL@{}}
\toprule
Loop & Generated layout \\
\midrule
\code{while (c) s} & \code{cond}$\star$: test \code{c}, if 0 jump to \code{end}; \code{s}; \code{j cond}; \code{end}: \\
\code{do s while (c);} & \code{body}: \code{s}; \code{cont}$\star$: test \code{c}, if 0 jump to \code{end}; \code{j body}; \code{end}: \\
\code{for (i; c; u) s} & \code{i}; \code{cond}: test \code{c}, if 0 jump to \code{end}; \code{s}; \code{cont}$\star$: \code{u}; \code{j cond}; \code{end}: \\
\bottomrule
\end{tabularx}
\end{center}

\snippetafter{int l_cond, l_body}
\snippet[language=yacc]{../parser.y}{case S_FOR:}{break;}

In a \code{for} loop, \code{continue} must jump to the step \code{u}, not to the condition; otherwise the step
would be skipped and the loop would never end. Each part may be omitted, and an omitted condition simply
emits no test, giving an infinite loop.

The targets of \code{break} and \code{continue} are kept on two stacks. A loop pushes its labels on entry and
pops them on exit, so \code{break} always reaches the innermost loop:

\snippet[language=yacc]{../parser.y}{static void push_loop}{@end}

\subsection{Switch}

\snippet[language=yacc]{../parser.y}{static void gen_switch}{@end}

A \code{switch} is emitted in two parts:

\begin{enumerate}
  \item \textbf{Dispatch.} The value is stored in its own frame slot (the extra slot reserved by
        \code{count_slots}). Each \code{case} value is evaluated, XOR-ed with the stored value, and a zero
        result jumps to that case's label. If nothing matches, control goes to \code{default}, or to the end
        when there is no \code{default}.
  \item \textbf{Bodies.} The case bodies are emitted in source order, each preceded by its label. Because they
        are contiguous, a case without \code{break} falls through into the next one, exactly as in C.
\end{enumerate}

The value lives in the frame rather than a register because evaluating a \code{case} value overwrites
\code{t0} and \code{t1}. A \code{switch} pushes only a \code{break} target, not a \code{continue} target, so a
\code{continue} inside a \code{switch} applies to the enclosing loop, again as in C.

\section{Functions}

\subsection{Calls}

\snippet[language=yacc]{../parser.y}{static void gen_call}{@end}

Arguments are evaluated left to right and pushed one by one. Only after all of them are evaluated are they
popped, in reverse order, into \code{a(n-1)} \ldots\ \code{a0}. Loading \code{a0}--\code{a7} directly would not
work, because a later argument may itself be a call --- \code{gcd(b, modulo(&a, &b))} --- and would overwrite
the registers already loaded. After the call, the result is copied from \code{a0} to \code{t0}, following the
``result in \code{t0}'' rule. Calls to the intrinsics are diverted to \code{gen_intrinsic} (Bonus).

\subsection{Recursion}

Recursion needs no special handling. Each call has its own frame with its own parameters and locals,
temporaries live on the stack, and \code{ra} and \code{s0} are saved by the prologue. The public test
Function/0 (\code{gcd}) and my test \file{calls} (\code{fact(10)}) both exercise recursion.

\subsection{Declarations}

The test programs declare prototypes before the definitions. \code{gen_program} adds every \code{D_FUNC} to
the symbol table (recording whether it returns a pointer) and emits code only for those with a body. A
function that is called without being declared, such as \code{digitalWrite} or \code{delay}, is not an error:
the generator emits \code{jal} and leaves the symbol to the linker.

\section{Global Variables}

\snippet[language=yacc]{../parser.y}{static void gen_global}{@end}

Globals live in \code{.data}, one \code{.dword} per element, with \code{.zero} filling elements that have no
initialiser. Initialisers must be integer constants, optionally negative; this is a limit of the assembler,
not an optimisation. Globals are read and written through \code{la t0, name}.

\section{Target Parameters}

\snippet[language=yacc]{../parser.y}{static int W = 8;}{+5}

Every size-dependent decision goes through these variables; the code never hard-codes 8 or \code{ld}.
With \code{--rv32}, \code{main} calls \code{use_rv32()}, which switches them to 4, 2, \code{lw}, \code{sw} and
\code{.word}. This is required by the Bonus, and it also lets the whole HW3 test suite run on RV32 as an extra
check.

\section{Testing}

\subsection{Procedure}

\code{make test} (script \file{tests/run.sh}) runs each test program in the Docker container:

\begin{lstlisting}[language=bash,numbers=none]
./codegen < test.c                                       # -> codegen.S
riscv64-unknown-elf-gcc -o prog spike-sample/main.c codegen.S
spike pk prog | grep '^Arduino'                          # -> output.txt
diff output.txt test.expected
\end{lstlisting}

\file{main.c} is the TA's Spike sample. It implements \code{digitalWrite} and \code{delay} by printing lines
such as \code{Arduino digitalWrite(26, 1);}, and calls \code{codegen()}. With \code{TARGET=rv32}, the script
uses the RV32 compiler and proxy kernel instead.

\subsection{Expected output}

The expected output of each public test is written by hand from the comments in the test program, as four
Arduino log lines. It is cross-checked with \file{tests/native_expected.sh}, which compiles the same C program
with the host gcc (\file{native_harness.h} supplies \code{digitalWrite}, \code{delay} and \code{main}) and runs
it; the native results match the hand-written files exactly. The expected output of my own tests comes from
this native run rather than from hand calculation.

\subsection{Public tests}

\begin{center}
\begin{tabularx}{\linewidth}{@{}lLl@{}}
\toprule
Test & Features & Expected log \\
\midrule
Basic/0 & function structure, Arduino calls, \code{HIGH}/\code{LOW} & (26,1) 1000 (26,0) 1000 \\
ArithmeticExpression/0, 1 & locals, arithmetic, parentheses & e.g.\ (26,1) 3000 (26,0) 3000 \\
Pointer/0, 1 & \code{&}, \code{*}, negative division, re-pointing & e.g.\ (28,1) 4000 (28,0) 4000 \\
Jump/0 & \code{for}, \code{while}, \code{if}/\code{else}, arrays, \code{*(arr+i)}, block-local declarations & (26,1) 3000 (26,0) 1000 \\
Jump/1 & \code{do-while}, nested \code{if}, \code{fib + i - 1} & (28,1) 6000 (28,0) 4000 \\
Function/0 & recursive \code{gcd}, \code{const int *} parameters, address of a parameter & (28,1) 6000 (28,0) 1000 \\
Function/1 & arrays as arguments, recursion & (26,1) 8000 (26,0) 3000 \\
\bottomrule
\end{tabularx}
\end{center}

Here (26,1) stands for \code{Arduino digitalWrite(26, 1);} and a bare number for \code{Arduino delay(...)}.

\subsection{My own tests}

The hidden tests are described as ``similar to level A (Function)''. To cover what the public tests do not,
I added six programs:

\begin{center}
\begin{tabularx}{\linewidth}{@{}lL@{}}
\toprule
Test & Coverage \\
\midrule
\file{ops} & every arithmetic, bitwise, comparison and logical operator; \code{/}, \code{\%} and \code{>>} on negative values \\
\file{control} & \code{continue}, \code{break}, nested \code{if}-\code{else}, \code{do-while} with \code{continue}, \code{switch} fall-through and \code{default}, \code{for(;;)} \\
\file{calls} & six parameters, calls as arguments, recursive \code{fact(10)}, swap through pointers, assigning through a returned pointer \\
\file{pointer_array} & pointer \code{++}/\code{--}, pointer difference, \code{&a[i]}, array initialisers, global scalars and arrays \\
\file{scope_incdec} & block-scope shadowing, values of prefix and postfix \code{++}/\code{--}, chained assignment \\
\file{short_circuit} & \code{&&} and \code{||} skip a right-hand call with a side effect \\
\bottomrule
\end{tabularx}
\end{center}

\section{Limitations}

\begin{center}
\begin{tabularx}{\linewidth}{@{}LL@{}}
\toprule
Limitation & Reason \\
\midrule
Only integers and single-level pointers; float and string literals are reported as errors & not used by the test programs \\
Arrays are one-dimensional with a constant length & not used by the test programs \\
At most 8 parameters & more would be passed on the stack; not used by the test programs \\
On RV64, \code{int} is computed in 8 bytes, so overflow differs from a 32-bit \code{int} & one word per value is a deliberate simplification; no test relies on overflow \\
The generated code is long & inherent in the stack-machine design; the specification forbids optimisation \\
\bottomrule
\end{tabularx}
\end{center}

\end{document}
